\section{Building the panel point-in-time}\label{sec:panel}

\begin{figure*}[tb]
\centering
\includegraphics[width=\textwidth]{fig8_vessel_year}
\caption{One carrier's \vesselFigureYear{} as the three pairings. Each row is a
terminal, US export terminals above the rule and European import terminals below.
A gold bar is a berth visit that ends with a laden departure, a black bar one that
ends in ballast, and a hatched bar the anchorage wait before berthing. Solid lines
are laden legs from a US departure to the next European zone entry; dashed lines are
the ballast return. This vessel made \vesselFigureVisits{} visits,
\vesselFigureLadenLegs{} laden and \vesselFigureBallastLegs{} ballast legs, from
\vesselFigureEvents{} events; the US events are from the NOAA archive and the
European ones from GFW, which is why the European visits carry no anchorage wait.}
\label{fig:vessel}
\end{figure*}

Figure~\ref{fig:vessel} shows the three pairings on one vessel's year. Every
signal is built from them twice. A \term{physical} basis uses hindsight and serves
validation. A \term{knowable} basis uses, on each day $d$, only what could have been
observed on day $d$, and serves modelling.

The headline signals are gas volumes. Let leg $i$ carry capacity $V_i$ (m$^3$ of
LNG), depart at $t^{\mathrm{dep}}_i$ and arrive at $t^{\mathrm{arr}}_i$. The
physical at-sea stock on day $d$ is
\begin{equation}
  S^{\mathrm{phys}}_d = \sum_i V_i \,
    \mathbf{1}\!\left[t^{\mathrm{dep}}_i \le d < t^{\mathrm{arr}}_i\right],
  \label{eq:stock-phys}
\end{equation}
banded by the arrival zone. The knowable stock replaces the arrival test with
what is known on $d$: a leg is at sea if it has departed and no arrival has been
\emph{observed} by $d$, and it is banded by destination only if a destination was
declared by $d$, otherwise it falls in the unresolved band. Loading and discharge
rates are amortised flows: a visit that berths for $H_j$ hours deposits $V_j/H_j$ per
hour across those hours, so that a closed visit integrates to exactly one cargo,
\begin{equation}
  L_d = \sum_j V_j\, \frac{\lvert [d, d+1) \cap [t^{\mathrm{moor}}_j, t^{\mathrm{dep}}_j) \rvert}{H_j},
  \label{eq:flow}
\end{equation}
and an open visit estimates $H_j$ from the terminal's closed-visit mean until it
departs, capped at one cargo. Durations (queue waits, voyage ages) are averaged
across terminals with weights equal to the number of legs or visits in each cell,
$\bar{x}_d = \sum_k n_k x_k / \sum_k n_k$; volumes are summed.

Two properties of voyage-derived features make the knowable basis difficult to
construct correctly. Both were discovered during the project rather than
anticipated, and Figure~\ref{fig:knowable} draws them.

\begin{figure*}[tb]
\centering
\resizebox{0.66\textwidth}{!}{% The knowable bound on one voyage leg, and the vintage hazard beneath it.
\begin{tikzpicture}[x=1cm, y=1cm, font=\scriptsize,
  arr/.style={-{Latex[length=1.5mm]}, draw=ink, line width=0.5pt}]
% ---- one leg on the event-time axis
\draw[arr] (0, 0) -- (8.3, 0) node[right, font=\tiny] {event time};
\foreach \x/\l in {0.9/$t_{\mathrm{dep}}$, 4.3/$d$, 6.7/$t_{\mathrm{arr}}$} {
  \draw[ink] (\x, -0.08) -- (\x, 0.08); \node[below, font=\tiny] at (\x, -0.1) {\l};
}
\draw[fill=palegold, draw=ink] (0.9, 1.05) rectangle (6.7, 1.45);
\node[left, font=\tiny] at (0.8, 1.25) {physical};
\node[font=\tiny] at (3.8, 1.25) {closed leg: laden, EU-bound, duration $t_{\mathrm{arr}}-t_{\mathrm{dep}}$};
\draw[fill=palegrey, draw=ink] (0.9, 0.4) rectangle (4.3, 0.8);
\draw[draw=ink!60, dashed] (4.3, 0.4) -- (6.7, 0.4) -- (6.7, 0.8) -- (4.3, 0.8);
\node[left, font=\tiny] at (0.8, 0.6) {knowable at $d$};
\node[font=\tiny] at (2.6, 0.6) {open leg, age $d-t_{\mathrm{dep}}$};
\node[font=\tiny, text=ink!70] at (5.5, 0.6) {destination unresolved};
\draw[dashed, draw=red] (4.3, -0.45) -- (4.3, 1.75);
\node[above, text=red, font=\tiny] at (4.3, 1.75) {as-of $d$: nothing to the right has happened yet};
% ---- the vintage axis
\draw[arr] (0, -1.35) -- (8.3, -1.35) node[right, font=\tiny] {vintage (years)};
\foreach \x/\l in {1.2/2018, 3.4/2020, 5.6/2024, 7.2/2026} {
  \draw[ink] (\x, -1.43) -- (\x, -1.27); \node[below, font=\tiny] at (\x, -1.45) {\l};
}
\draw[fill=paleblue, draw=ink] (7.2, -1.15) rectangle (8.3, -0.8);
\node[font=\tiny, left] at (7.1, -0.97) {destination-declaration feed exists};
\draw[dashed, draw=red] (3.4, -1.6) -- (3.4, -0.7);
\node[font=\tiny, text=red, above] at (3.4, -0.7) {at a 2020 as-of the table itself must not exist};
\end{tikzpicture}
}
\caption{Top: one leg on the event-time axis. In hindsight it is a closed, EU-bound
leg of known duration; at an as-of date $d$ before arrival it is an open leg of
known age and unknown destination, and the at-sea stock counts it as such. Bottom:
the vintage axis. A table that begins to exist in 2026 must return nothing at a
2020 as-of date even though every row in it is correctly dated.}
\label{fig:knowable}
\end{figure*}

\subsection{The unit of observation is defined in retrospect}

A voyage leg does not exist as a \emph{completed} unit until both a departure and
the next arrival have been observed. At time $t$, a vessel that departed five days
earlier is an open leg of known age and unknown duration and destination; that
open leg is itself a feature, and the at-sea stock of
Equation~\eqref{eq:stock-phys} counts it from the day it departed. What does not
yet exist is any statistic that conditions on completion: its duration, its arrival
zone, its floating-storage label. In hindsight the same vessel has a complete leg
with all three. Any statistic over legs therefore conditions on completion unless
the population is explicitly bounded to what was observable at $d$,
\begin{equation}
\begin{aligned}
  \mathcal{L}(d) &= \{\, i : t^{\mathrm{dep}}_i \le d \,\},\\
  t^{\mathrm{arr}}_i &\text{ attached only if } t^{\mathrm{arr}}_i \le d .
\end{aligned}
\label{eq:bound}
\end{equation}
Figure~\ref{fig:leakage} quantifies the effect. At an as-of date of 1~June~2020, a
loader that pairs all departures with all arrivals and then filters by date returns
\leakNaive{} legs. Of these, \leakFutureClosed{} (\leakFutureShare) are closed by
arrivals that had not yet occurred. With the population bounded by
Equation~\eqref{eq:bound}, \leakBounded{} legs remain, none of them closed by the
future. Status labels move with the bound: the number of legs classified as
floating storage fell from 95 to 9, because that label is assigned in hindsight. At
an as-of date after the data maximum the two loaders agree on every count, which
confirms that the bound is a restriction of the same query and not a different one.

\begin{figure}[t]
\centering
\includegraphics[width=0.95\linewidth]{fig2_leakage}
\caption{Voyage legs visible at an as-of date of 1~June~2020 under the naive loader
and under the point-in-time bound. Recomputed from the event tables when the figure
is generated.}
\label{fig:leakage}
\end{figure}

\subsection{Feature availability changes across vintages}

The same naive loader attached \leakDeclared{} destination declarations to the 2020
legs. Declarations come from a live AIS static-data feed that began in 2026, so no
declaration existed in 2020. This failure differs from the first. A date filter on
the observations is satisfied, because each declaration row carries a timestamp
within the vessel's current voyage, but the table being joined should not exist at
that vintage. The bound therefore reconstructs the declaration from the raw static
messages available at $d$, and at a 2020 as-of date it returns none.

\subsection{What the knowable basis bounds}\label{sec:bounds}

The knowable basis bounds on event time. Event timestamps are back-dated to the fix
at which a transition first qualified, so a departure is dated to the vessel's last
fix inside the berth polygon and not to the later fix that confirmed it had left.
The gap between the two is the feature's own publication lag. We measured it
directly. The first fix after a departure event arrives at a median lag of zero
hours, and more than 24 hours later for 0.3\% of departures on the NOAA feed
($n=1{,}500$) and 1.5\% on the live feed ($n=343$). On the decade panel the lag is
immaterial. On the live feed it is a genuine data limitation: a vessel that drops
out of terrestrial coverage as it leaves the berth is only back-dated once it
reappears, and a satellite feed would remove the gap. This is one of the reasons the
live acquisition system of Section~\ref{sec:live} spends most of its effort on
keeping vessels near terminals in a slot.

A separate log records the knowable value of every signal as it was printed on each
day of live operation, so that the bound can be checked rather than assumed.
Comparing the rebuilt knowable series with those prints on days old enough to have
settled, the check flags \vintagePostDay{} rows whose value moved after a next-day
print. All of them lie in the live tail and are consistent with late-arriving fixes
back-dating departures. A further \vintageSameDay{} rows differ from prints taken
during the day, before that day's fixes had finished arriving. That comparison
measures intraday timing and not leakage; it is reported but not scored. The live
tail lies outside the modelling sample used in Section~\ref{sec:partb}.

\subsection{The modelling grid}

Prices, storage and the tanker signals arrive on different calendars. They are
assembled onto one daily grid by forward-fill under a per-cadence staleness ceiling.
An inner join is never used: \HH{} spot prints on 69\% of calendar days, and an inner
join across the price legs would drop a third of the history without any indication
that it had done so. Each external series is shifted to the date on which it became
public (EIA weekly storage by six days, daily spot by one, AGSI+ by one). The tanker
signals are shifted by zero, because the knowable basis is already point-in-time and
a second shift would double-count. Degree days are not shifted. Their role is to be
partialled out, so the relevant quantity is the weather that occurred and not the
weather that was forecast. Figure~\ref{fig:signals} shows the eight signals that
enter Section~\ref{sec:partb}, on this grid.

\begin{figure*}[tb]
\centering
\includegraphics[width=0.95\textwidth]{fig18_signals}
\caption{The eight pre-registered tanker signals, daily on the knowable basis,
2018 to the coverage horizon. The upward drift in the volume signals is the growth
of US export capacity together with the capture gradient of
Figure~\ref{fig:capture}; the spread scan removes both through the persistence
term. The step at the very end of the EU-bound and unresolved in-transit panels is
the coverage seam.}
\label{fig:signals}
\end{figure*}

\subsection{The coverage seam}\label{sec:seam}

The two backfill feeds end on \noaaEnds{} (NOAA) and \gfwEnds{} (GFW). After those
dates the signals are carried by the live feed alone, which contributes
\nLiveEvents{} events against \nBackfillEvents{} from the backfills
(Figure~\ref{fig:coverage}). The EU-bound stock at sea falls from
\stockDecTwentyFive{} to \stockJanTwentySix{} million\,m$^3$ between December~2025
and January~2026, a \coverageDropShare{} drop with no physical cause. This was found
by plotting Figure~\ref{fig:spread} and was not caught by any automated check. The
primary modelling sample is therefore truncated at \noaaEnds. Full-panel results are
reported in Appendix~\ref{app:fullpanel} and change no conclusion.

\begin{figure*}[tb]
\centering
\includegraphics[width=0.95\textwidth]{fig6_coverage}
\caption{Port events per quarter by the feed that produced them, with the span of
each feed above. The NOAA and GFW backfills carry the decade; the live feed, split at
the cutover from a bounding-box subscription to a vessel-identifier filter, carries
only the tail. The dashed line is the coverage seam at which the primary sample ends.}
\label{fig:coverage}
\end{figure*}
